%%%%%%%%%%%%%%%%%%%%%%%%%%%%%%%%%%%%%%%%%%%%%%%%%%%%%%%%%%%%%%%%%%%%%%
% הגדרות משותפות למאגר ולדפי התרגילים שנבנים ממנו.
% נטען אחרי preamble/MotkePreamble, ודורש ש-\ifsol יוגדר לפני \documentclass:
%     \newif\ifsol
%     \soltrue      % או \solfalse
%%%%%%%%%%%%%%%%%%%%%%%%%%%%%%%%%%%%%%%%%%%%%%%%%%%%%%%%%%%%%%%%%%%%%%
\tcbuselibrary{theorems,breakable}
\usepackage{comment}
\usepackage{needspace}
\usetikzlibrary{calc}
\usetikzlibrary{arrows.meta}

% רמז לסעיף: \hint{...} בשורה נפרדת, לפני ה-solbox של הסעיף.
% במאגר (\soltrue) הרמז מודפס בכחול; בחוברת לסטודנטים (\solfalse) הוא מוסתר, ובאתר הוא נפתח בלחיצה.
\ifsol
\newcommand{\hint}[1]{\par{\small\color{blue!55!black}\textbf{רמז:}\enskip #1}\par\smallskip}
\else
\newcommand{\hint}[1]{}
\fi

\ifsol
\newtcolorbox{solbox}{
    enhanced,
    breakable,
    colback=gray!8,
    borderline east={2pt}{0pt}{gray!80},
    borderline west={0.3pt}{0pt}{gray!80},
    borderline north={0.3pt}{0pt}{gray!80},
    borderline south={0.3pt}{0pt}{gray!80},
    frame hidden,
    sharp corners,
    fontupper=\small,
    before upper={\textbf{פיתרון:}\enskip},
    grow to left by=0cm,
    grow to right by=0cm}
\else
\excludecomment{solbox}
\fi

\newcommand{\lvec}{\overrightarrow}
\renewcommand{\labelitemi}{$\diamond$}

% שם הקובץ של השאלה, בקצה השמאלי של השורה שמעליה (משמש ב-bank.tex).
% את השם הזה כותבים ב-\input כדי לשלב את השאלה בדף תרגילים.
\newcommand{\qtag}[1]{{\footnotesize\color{gray}\hfill קובץ: \textenglish{\ttfamily #1}}\par\nopagebreak[4]\smallskip}
\newcommand{\bankq}[2]{\item \qtag{#1}\input{#2}}

% \R{...} (טקסט עברי בתוך נוסחה) מגיע מ-babel ואינו מוגדר ב-MotkePreamble
\providecommand{\R}[1]{\RL{#1}}

% הפניה לשאלה אחרת במאגר לפי המזהה שלה, למשל \qref{int-013}
\newcommand{\qref}[1]{\textenglish{#1}}

% שאלה שהוצעה כדי לסגור פער בין מאגר המבחנים לדפי התרגילים (נכתבה ע"י Claude, טרם אושרה).
% כל הקובץ עטוף בסביבה הזו, ולכן השאלה מופיעה בתיבה ורודה בכל מקום שמשלבים אותה.
% לאחר אישור -- מוחקים את \begin{suggested} ואת \end{suggested}, ואת השורה "% suggested: yes" מכותרת הקובץ.
% עד אז השאלה מופיעה במאגר (bank.pdf) בלבד, ולא בחוברת לסטודנטים ולא באתר.
\definecolor{suggestpink}{RGB}{214,51,132}
\newtcolorbox{suggested}{
    enhanced,
    breakable,
    colback=suggestpink!7,
    colframe=suggestpink!70,
    boxrule=0.6pt,
    arc=2pt,
    left=6pt, right=6pt, top=4pt, bottom=4pt,
    before skip=2pt,
    before upper={{\color{suggestpink}\small\textbf{הצעה חדשה -- שאלה שנכתבה כדי לסגור פער מול מאגר המבחנים (טרם אושרה)}}\par\smallskip}}

% סימון לפתרון שעדיין לא נבדק ואושר (נכתב ע"י Claude). לאחר בדיקה -- מוחקים את השורה
% \unverified מתוך ה-solbox ואת השורה "% verified: no" מכותרת הקובץ.
\newcommand{\unverified}{{\color{red}\textbf{[טיוטה -- פתרון שטרם נבדק]}}\par}

% סעיפים בשתי עמודות -- רק בגרסה ללא פתרונות (בחוברת לסטודנטים).
% בגרסה עם פתרונות הסעיפים נשארים בעמודה אחת, כדי שהפתרונות לא יידחסו.
% שימוש: \begin{enumerate}\bMs{2} \item ... \eMs\end{enumerate}
\setlength{\multicolsep}{4pt plus 2pt minus 1pt}
\newcommand{\bMs}[1]{\ifsol\else\begin{multicols}{#1}\fi}
\newcommand{\eMs}{\ifsol\else\end{multicols}\fi}

%%%%%%%%%%%%%%%%%%%%%%%%%%%%%%%%%%%%%%%%%%%%%%%%%%%%%%%%%%%%%%%%%%%%%%
% פקודות של הקורס בהסתברות. אותן פקודות מוגדרות לאתר ב-site/macros.json.
%%%%%%%%%%%%%%%%%%%%%%%%%%%%%%%%%%%%%%%%%%%%%%%%%%%%%%%%%%%%%%%%%%%%%%
\providecommand{\N}{} \renewcommand{\N}{\mathbb{N}}
\providecommand{\set}{} \renewcommand{\set}[1]{\left\{#1\right\}}
\providecommand{\setof}{} \renewcommand{\setof}[2]{\left\{\,#1 \;\middle|\; #2\,\right\}}
\providecommand{\symdiff}{} \renewcommand{\symdiff}{\mathbin{\triangle}}
\providecommand{\comp}{} \renewcommand{\comp}[1]{\overline{#1}}
\providecommand{\Pow}{} \renewcommand{\Pow}{\mathcal{P}}
\renewcommand{\emptyset}{\varnothing}
\providecommand{\PP}{} \renewcommand{\PP}{\mathbb{P}}
\providecommand{\E}{} \renewcommand{\E}{\mathbb{E}}
\providecommand{\Var}{} \renewcommand{\Var}{\operatorname{Var}}
\providecommand{\Ber}{} \renewcommand{\Ber}{\mathrm{Ber}}
\providecommand{\Bin}{} \renewcommand{\Bin}{\mathrm{Bin}}
\providecommand{\Geo}{} \renewcommand{\Geo}{\mathrm{Geo}}
\providecommand{\Poi}{} \renewcommand{\Poi}{\mathrm{Poi}}
% התשובה הסופית של סעיף, במסגרת
\providecommand{\ans}{} \renewcommand{\ans}[1]{\boxed{#1}}

% דיאגרמת ון לשלוש קבוצות A, B, C. הארגומנט הראשון מונה את האזורים הצבועים: 1 = בתוך הקבוצה,
% 0 = מחוצה לה (לפי הסדר A/B/C). למשל \venn{1/1/0,0/1/1}{(ג)}. באתר הדיאגרמה מצוירת ב-SVG.
\definecolor{venncolor}{RGB}{35,72,168}
\newcommand{\vclip}[2]{\ifnum#2=1\relax\begin{scope}\clip #1;\else
  \begin{scope}[even odd rule]\clip (-1.6,-2.75) rectangle (3.0,1.65) #1;\fi}
\newcommand{\vennatom}[3]{\vclip{\vA}{#1}\vclip{\vB}{#2}\vclip{\vC}{#3}%
  \fill[venncolor!30] (-1.6,-2.75) rectangle (3.0,1.65);\end{scope}\end{scope}\end{scope}}
\newcommand{\venn}[2]{%
\begin{tikzpicture}[scale=0.42,baseline=(current bounding box.center)]
  \def\vA{(0,0) circle (1.1)}\def\vB{(1.4,0) circle (1.1)}\def\vC{(0.7,-1.15) circle (1.1)}
  \foreach \a/\b/\c in {#1} {\vennatom{\a}{\b}{\c}}
  \draw (-1.6,-2.75) rectangle (3.0,1.65);
  \draw[thick]\vA; \draw[thick]\vB; \draw[thick]\vC;
  \node at (-0.55,0.55) {\scriptsize $A$}; \node at (1.95,0.55) {\scriptsize $B$}; \node at (0.7,-1.85) {\scriptsize $C$};
  \node[below] at (0.7,-2.75) {\small #2};
\end{tikzpicture}}
